\documentclass[letterpaper,11pt]{article}

% T1: acentos como glifo único (í, no ı+´) y versalitas extraíbles como texto.
% Sin esto, un ATS o un Ctrl+F por "clínicos" no hace match.
\usepackage[T1]{fontenc}
\usepackage[empty]{fullpage}
\usepackage[spanish]{babel}
\usepackage{titlesec}
\usepackage[hidelinks]{hyperref}
\usepackage{enumitem}
\usepackage{tabularx}
\usepackage{charter}


% ==================== MÁRGENES ====================

\addtolength{\oddsidemargin}{-0.5in}
\addtolength{\evensidemargin}{-0.5in}
\addtolength{\textwidth}{1in}
\addtolength{\topmargin}{-0.5in}
\addtolength{\textheight}{1.3in}


% ==================== FORMATO DE SECCIONES ====================

% Mayúsculas reales en vez de \scshape: Charter no trae versalitas propias,
% las simula escalando la misma fuente y pdftotext lee el cambio de tamaño
% como espacio ("E XPERIENCIA"). \MakeUppercase extrae limpio.
\titleformat{\section}
  {\vspace{-4pt}\raggedright\large}
  {}
  {0em}
  {\MakeUppercase}
  [\titlerule \vspace{-5pt}]


% ==================== COMANDOS PROPIOS ====================

\newcommand{\resumeItem}[1]{%
  \item\small{#1 \vspace{-2pt}}%
}

\newcommand{\resumeSubheading}[4]{%
  \vspace{-2pt}\item
  \begin{tabular*}{0.97\textwidth}[t]{l@{\extracolsep{\fill}}r}
    \textbf{#1} & #2 \\
    \textit{\small#3} & \textit{\small #4} \\
  \end{tabular*}\vspace{-7pt}%
}

\newcommand{\resumeProjectHeading}[2]{%
  \item
  \begin{tabular*}{0.97\textwidth}{l@{\extracolsep{\fill}}r}
    \small#1 & #2 \\
  \end{tabular*}\vspace{-7pt}%
}

\newcommand{\resumeSubHeadingListStart}{%
  \begin{itemize}[leftmargin=0.15in, label={}]%
}
\newcommand{\resumeSubHeadingListEnd}{%
  \end{itemize}\vspace{-3pt}%
}

\newcommand{\resumeItemListStart}{\begin{itemize}}
\newcommand{\resumeItemListEnd}{\end{itemize}\vspace{-3pt}}

\raggedright
\raggedbottom
\setlength{\tabcolsep}{0in}


% ==================== DOCUMENTO ====================

\begin{document}


% -------------------- ENCABEZADO --------------------

\begin{center}
  \textbf{\Huge Felipe Farias Escobar} \\ \vspace{1pt}
  \small Ingeniero de Software Semi Senior \\ \vspace{1pt}
  Santiago, Chile \\
  \href{mailto:felipe.farias.e1@gmail.com}{felipe.farias.e1@gmail.com} $|$
  \href{https://www.linkedin.com/in/FelipeFa6/}{Linkedin.com/in/FelipeFa6} $|$
  \href{https://github.com/FelipeFa6}{Github.com/FelipeFa6}
\end{center}\vspace{-6pt}


% -------------------- RESUMEN PROFESIONAL --------------------

\section{Resumen Profesional}

\small{
  Ingeniero de Software con 4 años de experiencia construyendo sistemas
  backend desde cero en \textbf{Go} y \textbf{Python}, con soluciones en
  producción en los sectores de salud y energía.
}


% -------------------- EXPERIENCIA PROFESIONAL --------------------

\section{Experiencia Profesional}

\resumeSubHeadingListStart

  \resumeSubheading
    {Ingeniero de Software Semi Senior}{Junio 2025 -- Julio 2026}
    {INDRA}{Híbrido}
  \resumeItemListStart

    \resumeItem{Mantenimiento de microservicios en \textbf{Go} y tickets de
    alta variabilidad tecnológica (\textbf{Python}/Django,
    \textbf{TypeScript}/NestJS y React, \textbf{AWS}), con participación en
    \textbf{arquitectura}, \textbf{code review} y optimización de queries.
    Despliegues de microservicios sobre \textbf{GCP} (Compute Engine) con
    administración de máquinas por SSH.}

    \resumeItem{Implementación de la solución de \textbf{Calidad del Dato}
    junto al arquitecto de datos: reglas de validación, rutinas de
    verificación y listeners \textbf{event saga} (Java/Quarkus, Python).
    Calidad de datos de \textbf{75\% a 98\%} en plataformas del
    \textbf{Coordinador Eléctrico Nacional}, validada externamente.}

  \resumeItemListEnd

  \resumeSubheading
    {Ingeniero de Software}{Enero 2024 -- Junio 2025}
    {Instituto Nacional del Cáncer}{Presencial}
  \resumeItemListStart

    \resumeItem{Propuesta, diseño y construcción desde cero del
    \textbf{Integrador}, orquestador concurrente en \textbf{Go} que sincroniza
    \textbf{4 bases de datos Oracle} de sistemas clínicos aislados,
    procesando las solicitudes de la \textbf{API de un proveedor externo} a
    medida que llegan mediante endpoints HTTP propios con \textbf{Gin}.
    \textbf{Sigue en producción}.}

    \resumeItem{Patrón \textbf{event saga} con \textbf{compensación} ante
    fallos parciales, revirtiendo las bases ya actualizadas;
    \textbf{reintentos con encolado} persistente en base de datos y registro
    del estado de cada saga para \textbf{auditoría}. Concurrencia con
    \textbf{goroutines}; \textbf{tests unitarios y de integración} con
    \textbf{testify}.}

    \resumeItem{Habilitación de interoperabilidad \textbf{FHIR} con el
    Servicio de Salud Metropolitano Norte vía \textbf{Mirth Connect},
    eliminando la actualización manual de datos en producción.}

    \resumeItem{\textbf{CI/CD} con \textbf{Git Hooks} y migración a
    \textbf{Docker} (despliegue de \textbf{1 hora a 5 minutos});
    administración de servidores \textbf{Linux/Apache} y legacy \textbf{PHP
    5.7/CodeIgniter 3} migrado a \textbf{PHP 8/Laravel}.}

  \resumeItemListEnd

  \resumeSubheading
    {Desarrollador de Software Junior}{Julio 2022 -- Diciembre 2023}
    {Coolebra}{Remoto}
  \resumeItemListStart

    \resumeItem{Construcción desde cero de un \textbf{ETL} en \textbf{Python},
    \textbf{sin dependencias externas}, que auditaba \textbf{200 mil productos
    diarios} desde \textbf{PostgreSQL} hacia \textbf{OpenSearch} con historial
    de precios, en \textbf{2 horas} sobre \textbf{5 máquinas}
    (\textbf{Terraform}).}

    \resumeItem{Escalamiento del \textbf{scraping de más de 30 retailers}
    (BeautifulSoup, Selenium) con \textbf{ingeniería inversa} sobre
    \textbf{3 instancias EC2}, y despliegue de OpenSearch como servicio.}

    \resumeItem{Frontend en \textbf{React}: landing page y sección de
    productos con búsquedas contra OpenSearch; \textbf{Mercado Pago} para
    suscripciones.}

  \resumeItemListEnd

  \resumeSubheading
    {Tutor de Programación Orientada a Objetos III y Bases de Datos}
    {Noviembre 2021}
    {INACAP}{Presencial}
  \resumeItemListStart

    \resumeItem{Docencia en herencia y modelado orientado a objetos en
    \textbf{Python} y bases de datos relacionales.}

  \resumeItemListEnd

\resumeSubHeadingListEnd


% -------------------- PROYECTO ACTUAL --------------------

\section{Proyecto Actual}

\resumeSubHeadingListStart

  \resumeProjectHeading
    {\textbf{Petfi.io} $|$ Plataforma social de rescate y adopción}
    {Marzo 2026 -- Presente}
  \resumeItemListStart

    \resumeItem{Plataforma con más de \textbf{11 mil usuarios activos}.
    Aplicación web en \textbf{Flutter}, incluido el panel de administración y
    su sistema de roles, sobre \textbf{Firebase} (Firestore, Authentication,
    Storage, Messaging, Hosting) con reglas de seguridad e índices
    compuestos.}

  \resumeItemListEnd

\resumeSubHeadingListEnd


% -------------------- EDUCACIÓN Y CERTIFICACIONES --------------------

\section{Educación y Certificaciones}

\resumeSubHeadingListStart

  \resumeProjectHeading
    {\textbf{INACAP} $|$ Ingeniería en Informática $|$ Titulado 2023}
    {2019 -- 2022}

  \resumeProjectHeading
    {\textbf{University of Helsinki} $|$ Full Stack Open, curso principal}
    {2024}

  \resumeProjectHeading
    {\textbf{Oracle} $|$ OCI 2025 Certified Foundations Associate}
    {2026}

\resumeSubHeadingListEnd


% -------------------- TECNOLOGÍAS --------------------

\section{Tecnologías}

\small{
  \textbf{Lenguajes}: Go, Python, TypeScript, JavaScript, SQL \\
  \textbf{Backend y Datos}: Gin, NestJS, Django, FastAPI, Oracle/PL-SQL,
  PostgreSQL, OpenSearch, Firestore \\
  \textbf{DevOps}: AWS (EC2), GCP (Compute Engine, Firebase), Docker,
  Terraform, CI/CD, Git Hooks, Linux \\
  \textbf{Patrones}: Event Saga con compensación y reintentos, goroutines,
  microservicios, ETL, integración de sistemas, FHIR
}


\end{document}
